% Root and A full literal audits passed.  The110-case independent
% raw-scope/orbit/literal-bar gate is complete and source-hash linked.
\subsection{Eight states for the nonsquare seven-element source}

Write $M=1+(I_2\times J_3)$ and
$R_3=\{1,p,q\}$, where $pq=q$ and $qp=p$.
For an $(R_3,M)$-biset $S$, the left retractions $P,Q$ have
one common image $U$.  Its natural right $M$-graph $\Gamma$
has vertices $U$ and one edge $Ps\to Qs$ for each
$s\in S\setminus U$.  The collage has the actual homotopy model
\[
 X=\operatorname{hocofib}(\Gamma_{hM}\longrightarrow BM).
\]
Throughout this subsection maps between such models are the
actual graph and monoid maps, not hypothetical finite replacements.

\begin{lemma}[Split source quotients and full three-column loops]
\label{C:two-by-three-eight-state-reductions}
In the preceding model suppose $X\simeq\mathbb CP^2$.
Put $K=I_2\times J_3$ and $W=H_1(\Gamma;\mathbb Q)$.
Then no orbit of $SK$ is a singleton, and the projective-plane
class has a nonzero invariant edge functional
$\lambda:W\to\mathbb Q$.  The following cases are impossible:
\begin{enumerate}
\item the right action factors through either split quotient
$M\to1+J_3$ or $M\to P_5$;
\item $H_1(\Gamma e;\mathbb Q)$ is one-dimensional, generated by
a loop whose corresponding core outside-state orbit has all three
columns.
\end{enumerate}
\end{lemma}
\begin{proof}
The graph-component exact sequence identifies the kernel of
the edge-functional map with the image of the ordinary
component-act cone.  Its degree-two products vanish by
Lemma~\ref{C:nine-unbounded-act-cups}.  Thus a class with nonzero
square has $\lambda\ne0$.  A singleton core orbit supplies a
$K$-fixed graph vertex, after applying $P$ if necessary; that
vertex is also $M$-fixed and gives the usual section obstruction.
For every $e\in K$, invariance gives $\lambda=\lambda e$.
Consequently $We=H_1(\Gamma e;\mathbb Q)$ is nonzero.

For either split quotient, its monoid section gives an actual
collage retract, using the same graph with the factored action.
Restriction of the chosen second class retains the same nonzero
edge functional.  The retract's second-cohomology pullback
therefore spans the projective-plane line.  Over $1+J_3$, both
corner spaces have zero rational second cohomology, so the
relative suspension kills all degree-two products.  Over $P_5$
use Theorem~\ref{C:P5-graph-cofiber-cups}.  In either case
the pulled-back nonzero class has square zero.  This argument
does not infer finite homotopy type for the retract.

For the second assertion write $e=x_{00}$,
$a=x_{10}-e$, and $b_j=x_{0j}-e$ for $j=1,2$.
The exact left projective resolution already proved in
Lemma~\ref{C:two-by-three-seven-state-graph} is
\[
0\longrightarrow (Re)^2\xrightarrow{(b_1,b_2)}R(1-e)
 \xrightarrow{{-}\,a}Re\longrightarrow\mathbb Q
 \longrightarrow0,\qquad R=\mathbb QM.
\]
Here $We$ is one-dimensional, generated by the loop $c_0$.
The degree-two tensor map has basis images
$c_1-c_0,c_2-c_0$.  The three full-column core loops have
distinct literal edge supports, so these two cycles are
independent.  Hence $\operatorname{Tor}_2^R(W,\mathbb Q)=0$.
The component-module terms in the augmented-graph exact
sequence have Tor degrees three and four; both vanish by
the displayed length-two resolution.  Therefore
$H_4(X;\mathbb Q)=0$, a contradiction.
\end{proof}

\begin{lemma}[The unequal-partition eight-state control]
\label{C:two-by-three-eight-state-control}
Suppose $SK$ consists of three two-state orbits $A,B,C$.
Both left retractions send $B$ to $A$ and fix $A,C$;
their column partitions are respectively $001,001,010$.
There are two remaining transient states.  This normal form
contains no projective-plane collage.
\end{lemma}
\begin{proof}
If $P=Q$, the graph is a union of bouquets and strict loop
doubling excludes a rank-one nonzero square, using the
component-act zero-product theorem.  Otherwise label the
core by $A=\{0,1\}$, $B=\{2,3\}$, $C=\{4,5\}$.
The common restrictions of $P,Q$ are
$(0,1,0,1,4,5)$.  Each transient has an arbitrary ordered
row-orbit profile in $\{A,B,C\}^2$.  Their unordered
multisets give exactly $45$ right-action normal forms.
All transient retraction values are tried.  Idempotence,
commutation with every right actor, and equality of the
two images are checked literally.  There are $126$ unordered
pairs on these sorted profiles.  Quotienting only by the
actual transient-state stabilizer and interchange of $p,q$
retains $110$ cases.

The producer and complete records are
\path{code/a_eight_state_nonsquare_graphs.py} and
\path{code/outputs/a_eight_state_nonsquare_graphs.json}.
Exactly $100$ have nonzero first mod-two homology,
$4$ have the wrong second Betti number, and the other
$6$ have explicit square primitives for their nonboundary
second cocycles.

The independent checker
\path{code/d_verify_eight_state_nonsquare_graphs.py}
does not import the producer.  For each transient it
tries all $6^2$ first-column images in the two rows,
forces the remaining images by multiplication, and
checks all band equations.  It retains both transient
labels throughout, obtaining $81$ ordered right actions.
It then tries all $8^2$ retraction assignments and every
ordered distinct common-image pair, yielding $420$ pairs.
The complete expansions of the $110$ stored cases under
both transient permutations and both left-label permutations
are disjoint and equal exactly to that fresh raw set.
The sorted-profile counts independently recover $45$ and $126$.

The checker rebuilds each identity-first eighteen-arrow
category and uses the already separate literal bar verifier
from \path{code/b_verify_rightzero_six_isolated.py}.  It checks
$110220$ typed associativity equations, $81840$ boundary-square
equations, $3840$ cocycle equations and $23232$ square-primitive
equations.  No producer or matrix-library routine is imported.
The complete atomic record is
\path{code/outputs/d_eight_state_nonsquare_graphs_independently_verified.json}.
Thus every case in precisely this normal form is excluded.
\end{proof}

\begin{theorem}[The nonsquare source with eight right-zero states]
\label{C:two-by-three-eight-state-right-zero}
Let $M=1+(I_2\times J_3)$.  An $(R_3,M)$-collage with
at most eight states cannot have type $\mathbb CP^2$.
\end{theorem}
\begin{proof}
At most seven states are already excluded by
Theorem~\ref{C:seven-source-seven-state-right-zero}.
Assume therefore $|S|=8$ and use the nonzero $\lambda$
of Lemma~\ref{C:two-by-three-eight-state-reductions}.
Each orbit in $SK$ is a two- or three-state quotient of
$J_3$.  Let $o$ be the number of these orbits.

If $S=SK$, every band action is row-independent, so it
factors through the split row-collapse $M\to1+J_3$.
The same holds with one core orbit, even in the presence
of transients: both row images necessarily enter that
same orbit with its prescribed column.  These cases are
excluded by the first clause of the reduction lemma.
There are at most four core orbits, and four would use
all eight states.  Thus only $o=2,3$ with transients remain.

For $o=2$, the graph $\Gamma e$ has positive first
homology only when one orbit $A$ is a vertex orbit and
the other $B$ supplies a loop.  Both retractions send
$B$ to $A$, so the column partition of $B$ refines that
of $A$.  If $|B|=3$, the full-column loop clause applies.
If $|B|=2$, both orbits have the same two-part partition.
Every transient band image enters one of these two
orbits, so the entire right action factors through the
split column-collapse $M\to P_5$.  Again it is excluded.

Now let $o=3$.  The idempotent graph has three states.
Its nonzero first homology leaves only one vertex and
two loops, or two vertices and one loop.  In particular
the sole edge in the two-vertex case cannot be nonloop:
that would make $We=0$.  If every edge of $\Gamma$ is
a loop, simultaneous loop doubling and the component
argument finish at once.  We retain only graphs with
at least one nonloop edge.

Suppose there is one vertex orbit $A$ and two loop
orbits $B,C$.  If the core has six states, all three
two-state partitions agree, because the maps to $A$
are column-equivariant.  The split $P_5$ quotient applies.
If the core has seven states, $A$ has size two and one
of $B,C$ has size three.  There is exactly one transient.
It must be the nonloop edge joining the two vertices
of $A$; a transient vertex or loop would leave an
all-loop graph.  The resulting graph is connected,
and its first-homology module is literally
$\mathbb Q[B]\oplus\mathbb Q[C]$, with the transient
edge a bridge.  Its cycle coinvariants have dimension
two, one for each core orbit.  The connected graph
comparison then gives $\dim H_2(X;\mathbb Q)=2$,
contrary to the projective-plane rank.

Finally there are two vertex orbits $A,C$ and one loop
orbit $B$, mapped to $A$.  If the core has seven states
and $|B|=3$, the full-column loop clause applies.  Otherwise
$A,B$ both have size two and $C$ has size three.
There is one transient nonloop edge.  Its endpoints cannot
lie in different core vertex orbits: every band image of
the edge is a core state, on which $P=Q$, while distinct
core orbits remain distinct under the band.  If its
endpoints lie in $A$, all its band images lie in $A$ or
$B$; if they lie in $C$, all lie in $C$.  The graph
therefore has two nonempty component families preserved
separately by every actor.  The component act has at
least two orbits, giving $H_1(X;\mathbb Q)\ne0$.

If instead the core has six states, it has two transients.
The partitions of $A,B$ agree.  If that of $C$ also agrees,
the split $P_5$ quotient applies.  Otherwise relabeling the
three columns and the two states in each orbit gives
exactly $001,001,010$.  This is the complete normal form
of Lemma~\ref{C:two-by-three-eight-state-control}.
All core sizes and all surviving vertex/loop distributions
have now been excluded, proving the theorem.
\end{proof}
